
\documentclass{article}
\usepackage{colortbl}
\usepackage{makecell}
\usepackage{multirow}
\usepackage{supertabular}

\begin{document}

\newcounter{utterance}

\centering \large Interaction Transcript for game `guesswhat', experiment `Abs\_Level\_3', episode 8 with reka{-}flash{-}3.1{-}t0.0.
\vspace{24pt}

{ \footnotesize  \setcounter{utterance}{1}
\setlength{\tabcolsep}{0pt}
\begin{supertabular}{c@{$\;$}|p{.15\linewidth}@{}p{.15\linewidth}p{.15\linewidth}p{.15\linewidth}p{.15\linewidth}p{.15\linewidth}}
    \# & $\;$A & \multicolumn{4}{c}{Game Master} & $\;\:$B\\
    \hline

    \theutterance \stepcounter{utterance}  
    & & \multicolumn{4}{p{0.6\linewidth}}{
        \cellcolor[rgb]{0.9,0.9,0.9}{
            \makecell[{{p{\linewidth}}}]{
                \texttt{\tiny{[P1$\langle$GM]}}
                \texttt{Let's play a guessing game. Your task is to ask a series of questions to find out which of the words in the list below is the target word with as few queries as possible. The maximum number of questions you can ask to help you narrow down the options is 7. Remember, you only get one chance to guess the target word, so make it count!} \\
\\ 
\texttt{The rules of the game are: } \\
\\ 
\texttt{1. Only ask one question per turn. DO NOT ASK MORE THAN ONE QUESTION.} \\
\\ 
\texttt{2. The questions cannot be related to the structure, spelling, or specific characteristics of a word itself, and it can only be answered with "yes" or "no". } \\
\\ 
\texttt{3. Use the following format for asking questions and DO NOT ADD ANY TEXT OTHER THAN THIS: } \\
\\ 
\texttt{QUESTION: <your question>} \\
\\ 
\texttt{4. You have only one opportunity to guess the target word. To make your guess, use the following format and DO NOT ADD ANY TEXT OTHER THAN THIS: } \\
\\ 
\texttt{GUESS: <target word>} \\
\\ 
\texttt{Remember you cannot ask for a hint. Make your question brief. Let’s start!} \\
\\ 
\texttt{List of candidate words: {[}'sledding', 'cycling', 'football', 'boxing', 'swimming', 'surfing', 'canoeing', 'kayaking'{]}} \\
            }
        }
    }
    & \\ \\

    \theutterance \stepcounter{utterance}  
    & \multicolumn{4}{p{0.6\linewidth}}{
        \cellcolor[rgb]{0.9,0.9,0.9}{
            \makecell[{{p{\linewidth}}}]{
                \texttt{\tiny{[P1$\rangle$GM]}}
                \texttt{<reasoning>} \\
\texttt{Let's see. I need to figure out which word from the list is the target using up to 7 yes/no questions. The list is {[}'sledding', 'cycling', 'football', 'boxing','swimming','surfing', 'canoeing', 'kayaking'{]}.} \\
\\ 
\texttt{First, I should look for common characteristics among these words. All of them are activities or sports, I think. Let me check: sledding is a winter sport, cycling is a vehicle or activity, football is a team sport, boxing is combat, swimming is water{-}based, surfing is water, canoeing and kayaking are both paddling in water. So maybe categorizing them by type could help.} \\
\\ 
\texttt{Maybe start by splitting them into land vs water activities. Let's see: sledding (land?), cycling (land), football (land), boxing (land), swimming (water), surfing (water), canoeing (water), kayaking (water). Wait, sledding is on snow, so land. So if I ask if the target is a water{-}based activity, that could split the list into two groups. Let's count: water{-}based are swimming, surfing, canoeing, kayaking (4), land{-}based are the other four. If the answer is yes, then the target is among the four water activities. If no, then among the four land ones.} \\
\\ 
\texttt{But maybe there's a better way. Let's see. Alternatively, think about the equipment involved. For example, canoeing and kayaking both use paddles and are in water. Surfing uses a board, swimming uses a body. Boxing uses gloves, football a ball, cycling a bike, sledding a sled. But the rules say questions can't be about the structure or specific characteristics of the word itself. So I can't ask about equipment. So maybe focus on the activity type.} \\
\\ 
\texttt{Another approach: think about the number of people involved. Football is team, boxing is individual, swimming is individual, surfing individual, canoeing/kayaking can be individual or pairs. But again, can't ask about that. Maybe think about the environment: indoor vs outdoor? Football can be played both, but maybe some are more associated with outdoor. But not sure.} \\
\\ 
\texttt{Alternatively, think about the season. Sledding is winter, surfing is summer. But again, maybe not reliable.} \\
\\ 
\texttt{Wait, the rules say questions can't be about the structure} \\
            }
        }
    }
    & & \\ \\

    \theutterance \stepcounter{utterance}  
    & & & \multicolumn{2}{p{0.3\linewidth}}{
        \cellcolor[rgb]{0.9,0.9,0.9}{
            \makecell[{{p{\linewidth}}}]{
                \texttt{\tiny{[GM$|$GM]}}
                \texttt{Invalid format. Guesser must use the form 'QUESTION: ' or 'GUESS: '.} \\
            }
        }
    }
    & & \\ \\

    \theutterance \stepcounter{utterance}  
    & & & \multicolumn{2}{p{0.3\linewidth}}{
        \cellcolor[rgb]{0.9,0.9,0.9}{
            \makecell[{{p{\linewidth}}}]{
                \texttt{\tiny{[GM$|$GM]}}
                \texttt{abort game} \\
            }
        }
    }
    & & \\ \\

\end{supertabular}
}

\end{document}
